\providecommand{\rv}[1]{\text{#1}}

\section{Охрана труда}

\subsection{Производственная санитария, техника безопасности и пожарная профилактика}

Пользователи ПЭВМ могут подвергаться воздействию опасных и вредных производственных факторов, таких как повышенные уровни: электромагнитного, ультрафиолетового и инфракрасного излучения; статического электричества; повышенное или пониженное содержание аэроионов в воздухе рабочей зоны; повышенный или пониженный уровень освещённости рабочей зоны; содержание в воздухе рабочей зоны вредных веществ (оксид углерода, озон, аммиак, фенол, формальдегид); напряжение зрения, памяти, внимания; длительное статическое напряжение; монотонность труда; нерациональная организация рабочего места; эмоциональные перегрузки.

Требования к рабочим местам пользователей ПЭВМ приведены в Гигиеническом нормативе <<Показатели безопасности и безвредности факторов производственной среды и трудового процесса при работе с видеодисплейными терминалами и электронно-вычислительными машинами>> (постановление Совета Министров Республики Беларусь от 25.01.2021~г. №~37, в~ред. от 29.11.2022~г. №~829) [\ref{src:sanpin2}] и Типовой инструкцией по охране труда при использовании в работе офисного оборудования (постановление Министерства труда и социальной защиты от 14.04.2021~г. №~25) [\ref{src:office-instruction}].

Площадь помещения на одного пользователя ПЭВМ на базе плоских дискретных экранов (жидкокристаллических, плазменных) должна быть не менее 4{,}5~м\textsuperscript{2}.

\subsubsection{Метеоусловия}

В производственных помещениях, в которых работа с использованием ПЭВМ является основной (операторские, расчётные, посты управления, залы вычислительной техники), обеспечиваются оптимальные параметры микроклимата (таблица~\ref{tab:microclimate}).
Работа с компьютером относится к категории 1а (работы, производимые сидя и сопровождающиеся незначительным физическим напряжением, при которых расход энергии составляет до 120~ккал/ч, т.~е. до 139~Вт).

\begin{longtblr}[
  caption = {Оптимальные параметры микроклимата для помещений с ВДТ, ЭВМ и ПЭВМ},
  label = {tab:microclimate},
]{
  colspec = {|X[1,c,m]|X[1,c,m]|X[1.5,c,m]|X[1.5,c,m]|X[1.5,c,m]|},
  cells = {font=\small},
  width = \textwidth,
  hlines, vlines,
  rowhead = 1,
  row{1} = {halign=c},
}
Период года & Категория работ & Температура воздуха, \textdegree{}С, не более & Относительная влажность воздуха, \% & Скорость движения воздуха, м/с \\
Холодный & 1а & 22--24 & 40--60 & 0{,}1 \\
Тёплый & 1а & 23--25 & 40--60 & 0{,}1 \\
\end{longtblr}

Интенсивность теплового излучения работающих от нагретых поверхностей технологического оборудования, осветительных приборов, инсоляции на постоянных рабочих местах не превышает значений, указанных в таблице~\ref{tab:ir-radiation}.

\begin{longtblr}[
  caption = {Предельно допустимые уровни интенсивности излучения в инфракрасном и видимом диапазоне излучения на расстоянии 0{,}5~м со стороны экрана ВДТ, ЭВМ и ПЭВМ},
  label = {tab:ir-radiation},
]{
  colspec = {|X[1.4,l,m]|X[1,c,m]|X[1,c,m]|X[1,c,m]|},
  cells = {font=\small},
  width = \textwidth,
  hlines, vlines,
  rowhead = 1,
  row{1} = {halign=c},
}
Диапазоны длин волн & 400--760~нм & 760--1050~нм & свыше 1050~нм \\
Предельно допустимые уровни & 0{,}1~Вт/м\textsuperscript{2} & 0{,}05~Вт/м\textsuperscript{2} & 4{,}0~Вт/м\textsuperscript{2} \\
\end{longtblr}

Для создания нормальных метеорологических условий наиболее целесообразно уменьшить тепловыделения от самого источника~--~монитора, что предусматривается при разработке его конструкции.

В производственных помещениях для обеспечения необходимых показателей микроклимата предусмотрены системы отопления, вентиляции и кондиционирования воздуха.

\subsubsection{Вентиляция и отопление}

Воздух рабочей зоны помещения соответствует санитарно-гигиеническим требованиям по содержанию вредных веществ и частиц пыли, приведённым в ГОСТ~12.1.005--88 ССБТ <<Общие санитарно-гигиенические требования к воздуху рабочей зоны>> [\ref{src:gost-air}] и Гигиеническом нормативе <<Показатели безопасности и безвредности микроорганизмов продуцентов, микробных препаратов и их компонентов, вредных веществ в воздухе рабочей зоны и на кожных покровах работающих>> (постановление Совета Министров Республики Беларусь от 25.01.2021~г. №~37) [\ref{src:sanpin2}].

В помещениях проводится ежедневная влажная уборка и проветривание после каждого часа работы.

Уровни положительных и отрицательных аэроионов, а также коэффициент униполярности в воздухе всех помещений, где расположены ПЭВМ, соответствуют значениям, указанным в таблице~\ref{tab:aeroions}.

\begin{longtblr}[
  caption = {Аэроионный состав воздуха в зоне дыхания производственных и общественных помещений},
  label = {tab:aeroions},
]{
  colspec = {|X[1.3,l,m]|X[1.6,c,m]|X[1.1,c,m]|},
  cells = {font=\small},
  width = \textwidth,
  hlines, vlines,
  rowhead = 1,
  row{1} = {halign=c},
}
Уровни аэроионизации & Концентрация аэроионов, $\rho$ (ион/см\textsuperscript{3}) & Коэффициент униполярности, У \\
Минимально допустимые & $\rho\,\text{<<+>>} \geq 400 \quad \rho\,\text{<<--{}>>} \geq 600$ & \SetCell[r=2]{c,m} $0{,}4 < \rv{У} < 1{,}0$ \\
Максимально допустимые & $\rho\,\text{<<+>>} < 50\,000 \quad \rho\,\text{<<--{}>>} < 50\,000$ & \\
\end{longtblr}

Одним из мероприятий по оздоровлению воздушной среды является устройство систем вентиляции и отопления.
Задачей вентиляции является обеспечение чистоты воздуха и параметров метеорологических условий на рабочих местах.
Чистота воздушной среды достигается удалением загрязнённого или нагретого воздуха из помещения и подачей в него свежего воздуха.
Для поддержания нормального микроклимата необходим достаточный объём вентиляции, для чего предусматривается кондиционирование воздуха, осуществляющее поддержание постоянных параметров микроклимата в помещении независимо от наружных условий.

Параметры микроклимата поддерживаются в холодный период года за счёт системы водяного отопления с нагревом воды до 100~\textdegree{}С, а в тёплый~--~за счёт кондиционирования с параметрами, отвечающими требованиям СН~4.02.03--2019 [\ref{src:sn-heating}].

\subsubsection{Освещение}

В помещении при работе с ПЭВМ предусмотрены естественное и искусственное освещение.
Естественное освещение на рабочих местах осуществляется через световые проёмы, ориентированные преимущественно на север, северо-восток, восток, запад или северо-запад и обеспечивает коэффициент естественной освещённости не ниже 1{,}5~\%.
Оконные проёмы оборудованы регулируемыми устройствами типа жалюзи, занавесей.

Для внутренней отделки интерьера помещений используются материалы с коэффициентом отражения для потолка~--~0{,}7--0{,}8; для стен~--~0{,}5--0{,}6; для пола~--~0{,}3--0{,}5.

Искусственное освещение в помещениях осуществляется системой общего равномерного освещения, а при работе с документами применяется система комбинированного освещения.
Освещённость поверхности стола в зоне размещения документов составляет 300--500~люкс.
Освещённость поверхности экрана не более 300~люкс согласно СН~2.04.03--2020 [\ref{src:sn-light}].
В качестве источников света применяются люминесцентные лампы типа ЛБ.
Коэффициент запаса для осветительных установок общего освещения принимается равным 1{,}4, а коэффициент пульсации~--~не более 5~\%.

\subsubsection{Шум}

Источниками шума в помещениях, оборудованных ПЭВМ, являются принтеры, множительная техника и оборудование для кондиционирования воздуха, в самих ПЭВМ~--~вентиляторы систем охлаждения и трансформаторы.

В таблице~\ref{tab:noise} приведены допустимые уровни шума, которые обеспечиваются за счёт использования малошумного оборудования, применения звукопоглощающих материалов для облицовки помещений, а также различных звукопоглощающих устройств (перегородки и т.~п.).

\begin{longtblr}[
  caption = {Предельно допустимые уровни звука, эквивалентные уровни звука и уровни звукового давления в октавных полосах частот при работе с ВДТ, ЭВМ и ПЭВМ и периферийными устройствами},
  label = {tab:noise},
]{
  colspec = {|X[1.8,c,m]|X[0.7,c,m]|X[0.6,c,m]|X[0.6,c,m]|X[0.6,c,m]|X[0.6,c,m]|X[0.7,c,m]|X[0.7,c,m]|X[0.7,c,m]|X[0.7,c,m]|X[1.6,c,m]|},
  cells = {font=\footnotesize, cmd={\sloppy}},
  leftsep = 1pt,
  rightsep = 1pt,
  width = \textwidth,
  hlines, vlines,
  rowhead = 2,
  row{1,2} = {halign=c},
}
\SetCell[r=2]{c} Категория нормы шума & \SetCell[c=9]{c,m} {Уровни звукового давления, дБ,\\ в октавных полосах со среднегеометрическими частотами, Гц} & & & & & & & & & \SetCell[r=2]{c} Уровни звука и эквивалентные уровни звука, дБА \\
 & 31{,}5 & 63 & 125 & 250 & 500 & 1000 & 2000 & 4000 & 8000 & \\
I & 86 & 71 & 61 & 54 & 49 & 45 & 42 & 40 & 38 & 50 \\
II & 93 & 79 & 70 & 63 & 58 & 55 & 52 & 50 & 49 & 60 \\
III & 96 & 83 & 74 & 68 & 63 & 60 & 57 & 55 & 54 & 65 \\
IV & 103 & 91 & 83 & 77 & 73 & 70 & 68 & 66 & 64 & 75 \\
\end{longtblr}

\subsubsection{Электробезопасность}

Помещение вычислительного центра по степени опасности поражения электрическим током относится к помещениям без повышенной опасности.

Основные меры защиты от поражения током~--~изоляция и недоступность токоведущих частей; защитное заземление ($R_{\rv{з}} = 4$~Ом, ГОСТ~12.1.030--81) [\ref{src:gost-grounding}].

Первая помощь при поражениях электрическим током состоит из двух этапов: освобождение пострадавшего от действия тока и оказание ему доврачебной медицинской помощи.
После освобождения пострадавшего от действия электрического тока необходимо оценить его состояние.
Во всех случаях поражения электрическим током необходимо вызвать врача независимо от состояния пострадавшего.

\subsubsection{Излучение}

При работе с дисплеем ПЭВМ могут возникать следующие опасные факторы: электромагнитные и электростатические поля, ультрафиолетовое и инфракрасное излучение.

Уровни электромагнитных и электростатических полей на рабочих местах пользователей, создаваемые экранами ПЭВМ, не превышают предельно допустимых значений (таблица~\ref{tab:emf-screen}).

\begin{longtblr}[
  caption = {Предельно допустимые уровни электромагнитных полей от экранов ВДТ, ЭВМ и ПЭВМ},
  label = {tab:emf-screen},
]{
  colspec = {|X[2,l,m]|X[1.2,c,m]|},
  cells = {font=\small},
  width = \textwidth,
  hline{1,2,5,8,9} = {solid},
  vlines,
  rowhead = 1,
  row{1} = {halign=c},
}
Наименование параметра & Предельно допустимые уровни \\
Напряжённость электрического поля в диапазоне частот: & \\
\hfill 5~Гц--2~кГц\hfill\null & не более 25{,}0~В/м \\
\hfill 2--400~кГц\hfill\null & не более 2{,}5~В/м \\
Плотность магнитного потока магнитного поля в диапазоне частот: & \\
\hfill 5~Гц--2~кГц\hfill\null & не более 250~нТл \\
\hfill 2--400~кГц\hfill\null & не более 25~нТл \\
Напряжённость электростатического поля & не более 15~кВ/м \\
\end{longtblr}

Наиболее эффективным и часто применяемым методом защиты от электромагнитных излучений является установка экранов.
Экранируют либо источник излучения, либо рабочее место.
Часто экран устанавливают непосредственно на монитор.

При работе монитора на экране кинескопа накапливается электростатический заряд, создающий электростатическое поле.
При этом персонал, работающий с монитором, приобретает электростатический потенциал.
Заметный вклад в общее электростатическое поле вносят электризующиеся от трения поверхности клавиатуры и мыши; предельно допустимые уровни электромагнитных полей от периферийных устройств приведены в таблице~\ref{tab:emf-peripheral}.

\begin{longtblr}[
  caption = {Предельно допустимые уровни электромагнитных полей при работе с ВДТ, ЭВМ, ПЭВМ от клавиатуры, системного блока, манипулятора <<мышь>>, беспроводных систем передачи информации и иных периферийных устройств},
  label = {tab:emf-peripheral},
]{
  colspec = {|X[1.6,l,m]|X[1,c,m]|X[1,c,m]|X[1,c,m]|X[1,c,m]|X[1.2,c,m]|},
  cells = {font=\small},
  leftsep = 3pt,
  rightsep = 3pt,
  width = \textwidth,
  hlines, vlines,
  rowhead = 1,
  row{1} = {halign=c},
}
Диапазоны частот & 0{,}3--300~кГц & 0{,}3--3~МГц & 3--30~МГц & 30--300~МГц & 0{,}3--300~ГГц \\
Предельно допустимые уровни & 25~В/м & 15~В/м & 10~В/м & 3~В/м & 10~мкВт/см\textsuperscript{2} \\
\end{longtblr}

Уровни ультрафиолетового излучения, создаваемого экраном ВДТ, ЭВМ и ПЭВМ, не превышают предельно допустимых значений (таблица~\ref{tab:uv}).

\begin{longtblr}[
  caption = {Предельно допустимые уровни интенсивности излучения в ультрафиолетовом диапазоне на расстоянии 0{,}5~м со стороны экрана ВДТ, ЭВМ и ПЭВМ},
  label = {tab:uv},
]{
  colspec = {|X[1.4,l,m]|X[1,c,m]|X[1,c,m]|X[1,c,m]|},
  cells = {font=\small},
  width = \textwidth,
  hlines, vlines,
  rowhead = 1,
  row{1} = {halign=c},
}
Диапазоны длин волн & 200--280~нм & 280--315~нм & 315--400~нм \\
Предельно допустимые уровни & не допускается & 0{,}0001~Вт/м\textsuperscript{2} & 0{,}1~Вт/м\textsuperscript{2} \\
\end{longtblr}

\subsubsection{Пожарная безопасность}

По взрывопожарной и пожарной опасности помещения и здания для ЭВМ относятся к категории~Д согласно ТКП~474--2013 [\ref{src:tkp474}], а по степени огнестойкости~--~ко второй степени согласно СН~2.02.05--2020 [\ref{src:sn-fire}].

Для предотвращения распространения огня во время пожара с одной части здания на другую устраивают противопожарные преграды в виде стен, перегородок, дверей, окон.
Особое требование предъявляется к устройству и размещению кабельных коммуникаций.

Эвакуация персонала вычислительного центра осуществляется через эвакуационные выходы.
Количество и общая ширина эвакуационных выходов определяются в зависимости от максимального возможного числа эвакуирующихся через них людей и предельно допустимого расстояния от наиболее удалённого места возможного пребывания людей до ближайшего эвакуационного выхода согласно СН~2.03.05--2020 [\ref{src:sn-evacuation}].

Расчётное время эвакуации устанавливается по реальному расчёту времени движения одного или нескольких потоков людей через эвакуационные выходы из наиболее удалённых мест размещения людей.
Необходимое время эвакуации устанавливается на основе данных о критической продолжительности пожара с учётом степени огнестойкости здания, категории производства по взрывной и пожарной опасности.

Для ликвидации пожаров в начальной стадии применяются первичные средства пожаротушения: внутренние пожарные водопроводы, огнетушители типа ОВП-10, ОУ-2, асбестовые одеяла и др.
Нормы первичных средств пожаротушения для вычислительных центров приведены в таблице~\ref{tab:fire-extinguishers}.

\begin{longtblr}[
  caption = {Примерные нормы первичных средств пожаротушения для вычислительного центра},
  label = {tab:fire-extinguishers},
]{
  colspec = {|X[1.5,l,m]|X[1,c,m]|X[1.4,c,m]|X[1.2,c,m]|},
  cells = {font=\small},
  width = \textwidth,
  hlines, vlines,
  rowhead = 1,
  row{1} = {halign=c},
}
Помещение & Площадь, м\textsuperscript{2} & Углекислотные огнетушители ручные & Порошковые огнетушители \\
Вычислительный центр & 100 & 1 & 1 \\
\end{longtblr}

\subsection{Факторы, влияющие на исход поражения электрическим током}

Исход воздействия электрического тока зависит от следующих факторов: величины тока, длительности протекания электрического тока через тело человека, электрического сопротивления тела человека, рода и частоты тока, пути тока в организме человека.

Электрическое сопротивление тела человека определяется сопротивлением кожи и внутренних тканей.
Поверхностный слой кожи, называемый эпидермисом, состоящий в основном из мёртвых ороговевших клеток, обладает большим сопротивлением, которое и определяет общее сопротивление тела человека.
Сопротивление нижних слоёв кожи и внутренних тканей человека незначительно.
При сухой, чистой и неповреждённой коже сопротивление тела человека колеблется в пределах 2~тыс.--2~млн~Ом.
При увлажнении, загрязнении и при повреждении кожи сопротивление тела оказывается равным около 500~Ом (сопротивление внутренних тканей тела).
В расчётах сопротивление тела человека принимается равным 1000~Ом.

Величина тока, протекающего через тело человека, является главным фактором, от которого зависит исход поражения.
Ощутимый ток~--~человек начинает ощущать протекающий через него ток промышленной частоты 50~Гц относительно малого значения: 0{,}5--1{,}5~мА.
Неотпускающий ток~--~ток 10--15~мА вызывает сильные и весьма болезненные судороги мышц рук, которые человек самостоятельно преодолеть не в состоянии и оказывается как бы прикованным к токоведущей части.
При 25--50~мА действие тока распространяется и на мышцы грудной клетки, что приводит к затруднению и даже прекращению дыхания.
При длительном воздействии этого тока~--~в течение нескольких минут~--~может наступить смерть вследствие прекращения работы лёгких.
Фибрилляционный ток~--~при 100~мА ток оказывает непосредственное влияние и на мышцу сердца; при длительности протекания более 0{,}5~с такой ток может вызвать остановку или фибрилляцию сердца, т.~е. быстрые хаотические и разновременные сокращения волокон сердечной мышцы (фибрилл), при которых сердце перестаёт работать как насос.
В результате в организме прекращается кровообращение и наступает смерть.

Длительность протекания тока через тело человека влияет на исход поражения вследствие того, что со временем резко повышается ток за счёт уменьшения сопротивления тела.
Кроме того, длительное прохождение переменного тока нарушает ритм сердечной деятельности, вызывая трепетание желудочков сердца в связи с поражением нервов сердечной мышцы.

Род и частота тока в значительной степени определяют исход поражения.
Наиболее опасным является переменный ток с частотой 20--100~Гц.
При частоте меньше 20 или больше 100~Гц опасность поражения током заметно снижается.
Токи частотой свыше 500\,000~Гц не оказывают раздражающего действия на ткани и поэтому не вызывают электрического удара.
Однако они могут вызвать термические ожоги.
При постоянном токе пороговый ощутимый ток повышается до 6--7~мА, пороговый неотпускающий ток~--~до 50--70~мА, а фибрилляционный при длительности воздействия более 0{,}5~с~--~до 300~мА.

Путь прохождения тока через тело человека также влияет на исход поражения.
Наибольшую опасность представляет прохождение тока через жизненно важные органы (сердце, спинной мозг, органы дыхания и т.~д.) по пути <<рука--рука>> и <<рука--ноги>>, при этом ток проходит по кровеносным и лимфатическим сосудам, оболочкам нервных стволов и т.~д.
Менее опасен путь тока <<нога--нога>>.
